\documentclass[11pt,a4paper]{article}
\usepackage{fontspec}
\setmainfont{PT Serif}[Ligatures=TeX]
\setsansfont{Helvetica Neue}[Ligatures=TeX]
\setmonofont{Menlo}
\usepackage{polyglossia}
\setdefaultlanguage{russian}
\setotherlanguage{english}
\newfontfamily\cyrillicfont{PT Serif}[Ligatures=TeX]
\newfontfamily\cyrillicfontsf{Helvetica Neue}
\newfontfamily\cyrillicfonttt{Menlo}
\usepackage{amsmath,amssymb,amsthm}
\usepackage{geometry}
\geometry{margin=22mm}
\usepackage{booktabs}
\usepackage{enumitem}
\setlist{nosep,leftmargin=*}
\usepackage{microtype}
\usepackage{hyperref}
\hypersetup{colorlinks,linkcolor=black,urlcolor=blue!50!black}
\renewcommand{\arraystretch}{1.12}
\emergencystretch=2em
\newcommand{\E}{\mathbf{E}}
\newcommand{\Prob}{\mathbf{P}}
\newcommand{\D}{\mathbf{D}}
\newcommand{\Ind}{\mathrm{I}}
\newcommand{\Bern}{\mathrm{Bern}}
\newcommand{\Bin}{\mathrm{Bin}}
\DeclareMathOperator*{\argmax}{arg\,max}
\theoremstyle{plain}
\newtheorem{theorem}{Теорема}
\newtheorem{lemma}{Лемма}
\newtheorem{proposition}{Утверждение}
\theoremstyle{definition}
\newtheorem{remark}{Замечание}

\title{Оценка параметров одного LLM-теста:\\ функция правдоподобия, ОМП и EM-алгоритм}
\author{к задаче «Экономия на эвалах»}
\date{25 октября 2026}

\begin{document}
\sloppy
\maketitle

\paragraph{Об обозначениях.} Используются обозначения, принятые в курсе математической статистики мехмата: выборка $X=(X_1,\dots,X_N)$, параметрическое семейство $\{\Prob_\theta,\ \theta\in\Theta\}$, функция правдоподобия $L(\theta;X)$, оценка максимального правдоподобия $\hat\theta$, математическое ожидание и дисперсия $\E_\theta$, $\D_\theta$, индикатор $\Ind\{A\}$, биномиальное распределение $\Bin(n,s)$, распределение Бернулли $\Bern(p)$. Конспекты кафедры на момент написания были недоступны (teach-in.ru и math.msu.ru не отвечали), поэтому сверить обозначения буквально не удалось.

\section{Модель для одного теста}

Фиксируем тест $j$ и далее индекс $j$ опускаем. Релизов $N=70$, на каждом тест прогнан $r=5$ раз. Скрытая величина $I_i\in\{0,1\}$ --- есть ли в релизе $i$ проблема в сценарии теста; наблюдаемая $K_i\in\{0,\dots,r\}$ --- число красных прогонов. Модель:
\[
I_1,\dots,I_N\ \text{независимы},\quad I_i\sim\Bern(p);\qquad K_i\mid I_i=1\ \sim\ \Bin(r,s),\qquad K_i\mid I_i=0\ \equiv 0 .
\]
Параметр $\theta=(p,s)\in\Theta=[0,1]^2$. Ложных красных нет: при $I_i=0$ все прогоны зелёные.

\begin{remark}
В данных записан порядок прогонов внутри релиза, но при $I_i=1$ прогоны независимы и одинаково распределены, поэтому вектор исходов пяти прогонов зависит от $\theta$ только через $K_i$: число красных --- достаточная статистика для одного релиза. Правдоподобие по исходным данным отличается от правдоподобия по $K_i$ множителем, не зависящим от $\theta$.
\end{remark}

Распределение одного наблюдения:
\begin{equation}\label{eq:pmf}
\Prob_\theta(K=0)=1-p+p(1-s)^r,\qquad
\Prob_\theta(K=k)=p\binom rk s^k(1-s)^{r-k},\quad k=1,\dots,r .
\end{equation}
Нуль возникает двумя путями: проблемы нет, либо она есть, но ни разу не проявилась.

\section{Функция правдоподобия и ОМП}

Выборка $X=(K_1,\dots,K_N)$. Обозначим
\[
n_+=\sum_{i=1}^N\Ind\{K_i\ge1\},\qquad n_0=N-n_+,\qquad S=\sum_{i=1}^NK_i .
\]
По \eqref{eq:pmf} функция правдоподобия равна
\begin{align}
L(\theta;X)&=\bigl[1-p+p(1-s)^r\bigr]^{n_0}\prod_{i:\,K_i\ge1}p\binom r{K_i}s^{K_i}(1-s)^{r-K_i}\notag\\
&=\bigl[1-p+p(1-s)^r\bigr]^{n_0}\,p^{\,n_+}\,s^{S}(1-s)^{rn_+-S}\prod_{i:\,K_i\ge1}\binom r{K_i},\label{eq:lik}
\end{align}
а логарифмическая функция правдоподобия, с точностью до слагаемого, не зависящего от $\theta$, ---
\begin{equation}\label{eq:loglik}
\ell(p,s)=n_0\log\bigl[1-p+p(1-s)^r\bigr]+n_+\log p+S\log s+(rn_+-S)\log(1-s).
\end{equation}
По теореме факторизации $(n_+,S)$ --- достаточная статистика: сколько релизов с красным и сколько всего красных.

\begin{lemma}[разделение параметров]\label{lem:split}
Положим $q=\Prob_\theta(K\ge1)=p\bigl(1-(1-s)^r\bigr)$. Отображение $(p,s)\mapsto(q,s)$ взаимно однозначно на $(0,1]\times(0,1]$, и
\begin{equation}\label{eq:split}
\ell=\underbrace{n_+\log q+n_0\log(1-q)}_{\ell_1(q)}\ +\ \underbrace{S\log s+(rn_+-S)\log(1-s)-n_+\log\bigl(1-(1-s)^r\bigr)}_{\ell_2(s)} .
\end{equation}
\end{lemma}

\begin{proof}
$1-q=1-p+p(1-s)^r$, а $p=q/(1-(1-s)^r)$; подставляем в \eqref{eq:loglik}.
\end{proof}

Второе слагаемое $\ell_2(s)$ --- логарифм правдоподобия выборки $\{K_i:\ K_i\ge1\}$ из \emph{усечённого в нуле} биномиального распределения $\Prob(K=k\mid K\ge1)=\binom rk s^k(1-s)^{r-k}/(1-(1-s)^r)$.

\begin{proposition}[ОМП в явном виде]\label{prop:mle}
Пусть $n_+\ge1$ и $\bar K_+=S/n_+$ --- среднее число красных по релизам с красным. Тогда
\begin{equation}\label{eq:mle}
\hat q=\frac{n_+}N,\qquad
\hat s:\quad \frac{r\hat s}{1-(1-\hat s)^r}=\bar K_+,\qquad
\hat p=\min\Bigl\{\frac{\hat q}{1-(1-\hat s)^r},\,1\Bigr\}.
\end{equation}
Уравнение для $\hat s$ имеет единственный корень на $(0,1)$ при $1<\bar K_+<r$. При $\bar K_+=1$ оценка вырождается: $\hat s\to0$, $\hat p=1$; при $\bar K_+=r$ имеем $\hat s=1$, $\hat p=\hat q$. При $n_+=0$: $\hat p=0$, а $s$ не оценивается.
\end{proposition}

\begin{proof}
По лемме~\ref{lem:split} максимизация распадается. $\ell_1$ --- биномиальное правдоподобие, максимум в $\hat q=n_+/N$. Для $\ell_2$: $\partial\ell_2/\partial s=0$ равносильно $\E(K\mid K\ge1)=\bar K_+$, то есть уравнению \eqref{eq:mle}; функция $s\mapsto rs/(1-(1-s)^r)$ непрерывна, строго возрастает и пробегает $(1,r)$ при $s\in(0,1)$ (это среднее усечённого бинома, монотонное по $s$ как среднее экспоненциального семейства), откуда единственность. Возвращаясь к $p$, получаем $\hat p=\hat q/(1-(1-\hat s)^r)$; если это число больше единицы, максимум по $p$ на $[0,1]$ достигается на границе $p=1$.
\end{proof}

\begin{remark}
Из \eqref{eq:mle} следует $\hat p\hat s=S/(rN)$: ОМП точно воспроизводит выборочное среднее $\bar K=S/N$, потому что $\E_\theta K=rps$. Второе условие --- совпадение доли релизов с красным: $\hat p\,(1-(1-\hat s)^r)=n_+/N$. Эти два равенства --- те же уравнения, что у метода моментов по статистикам $\Ind\{K\ge1\}$ и $K$.
\end{remark}

\section{EM-алгоритм}

Полные данные --- пары $(I_i,K_i)$. Их логарифмическое правдоподобие
\begin{equation}\label{eq:complete}
\ell_{\rm c}(\theta)=\sum_{i=1}^N\Bigl[I_i\bigl(\log p+K_i\log s+(r-K_i)\log(1-s)\bigr)+(1-I_i)\log(1-p)\Bigr]+\mathrm{const}
\end{equation}
линейно по $I_i$, поэтому E-шаг сводится к подстановке условных математических ожиданий $w_i=\E_{\theta}(I_i\mid K_i)=\Prob_\theta(I_i=1\mid K_i)$.

\paragraph{E-шаг.} При $K_i\ge1$ проблема есть наверняка: $w_i=1$. При $K_i=0$ по формуле Байеса
\begin{equation}\label{eq:estep}
w_i=w_0(\theta)=\frac{p(1-s)^r}{1-p+p(1-s)^r}.
\end{equation}

\paragraph{M-шаг.} Максимизируем $Q(\theta\mid\theta^{(t)})=\E_{\theta^{(t)}}\bigl(\ell_{\rm c}(\theta)\mid X\bigr)$, то есть \eqref{eq:complete} с $I_i$, заменёнными на $w_i$. Обозначим $W=\sum_iw_i=n_+ + n_0w_0(\theta^{(t)})$ --- «эффективное» число проблемных релизов. Тогда
\begin{equation}\label{eq:mstep}
p^{(t+1)}=\frac WN,\qquad s^{(t+1)}=\frac{\sum_iw_iK_i}{rW}=\frac S{rW}.
\end{equation}
Первое --- средний вес, второе --- доля красных среди прогонов на релизах, взвешенных по $w_i$; в числителе $\sum_iw_iK_i=S$, так как $K_i=0$ там, где $w_i<1$.

\paragraph{Алгоритм.} Начальное приближение $p^{(0)}=n_+/N$, $s^{(0)}=S/(rn_+)$ (релизы без красного считаем чистыми); затем повторять \eqref{eq:estep}--\eqref{eq:mstep} до стабилизации.

\begin{theorem}
Последовательность $\theta^{(t)}$ не убывает по правдоподобию: $\ell(\theta^{(t+1)})\ge\ell(\theta^{(t)})$. Любая её предельная точка с $n_+\ge1$ и $1<\bar K_+<r$ есть ОМП \eqref{eq:mle}.
\end{theorem}

\begin{proof}
Монотонность --- общее свойство EM (Dempster, Laird, Rubin, 1977): $\ell(\theta)=Q(\theta\mid\theta^{(t)})-H(\theta\mid\theta^{(t)})$, где $H(\theta\mid\theta^{(t)})=\E_{\theta^{(t)}}\bigl(\log\Prob_\theta(I\mid X)\mid X\bigr)$, а $H(\theta\mid\theta^{(t)})\le H(\theta^{(t)}\mid\theta^{(t)})$ по неравенству Йенсена; M-шаг увеличивает $Q$, значит, и $\ell$.

Пусть $\theta=(p,s)$ --- неподвижная точка \eqref{eq:mstep}: $p=W/N$, $s=S/(rW)$, где $W=n_++n_0w_0(\theta)$. Перемножая, получаем $ps=S/(rN)$. Далее, $W-n_+=n_0w_0$ и $N-W=n_0(1-w_0)$, поэтому
\[
\frac{W-n_+}{N-W}=\frac{w_0}{1-w_0}=\frac{p(1-s)^r}{1-p}.
\]
Подставим $W=Np$: $N-W=N(1-p)$, и равенство даёт $Np-n_+=Np(1-s)^r$, то есть $p\bigl(1-(1-s)^r\bigr)=n_+/N$. Итак, неподвижная точка решает систему $ps=S/(rN)$, $p(1-(1-s)^r)=n_+/N$ --- в точности уравнения ОМП из замечания после утверждения~\ref{prop:mle}; при $1<\bar K_+<r$ их решение единственно, а по монотонности правдоподобие сходится к его максимуму. Если же $\bar K_+=1$, EM уходит к границе $s\to0$, $p\to1$, как и ОМП.
\end{proof}

\begin{remark}[зачем EM, если есть явная формула]
Для одного теста EM не нужен: ответ даёт утверждение~\ref{prop:mle}. EM полезен, когда параметры связаны между тестами --- например, общая чувствительность $s$ для группы тестов. Тогда E-шаг по каждому тесту тот же, а M-шаг объединяет: $p_j=W_j/N$, $s=\sum_jS_j\big/\bigl(r\sum_jW_j\bigr)$; явной формулы здесь уже нет, а EM остаётся монотонным.
\end{remark}

\section{Что получается на данных задачи}

В опубликованной истории ($N=70$, $M=40$ тестов) у типичного теста $n_+=2$ релиза с красным. У $9$ тестов $n_+=0$: для них $\hat p=0$, а $s$ не определено. Вырожденный случай $\bar K_+=1$ (все релизы с красным дали ровно по одному красному) на опубликованной истории не встретился, но в истории того же мира хотя бы один такой тест есть в $78\%$ случаев, в среднем $1.5$ теста на историю: для него ОМП даёт $\hat s=0$, $\hat p=1$ --- тест выглядит как «ломается всегда, но почти не виден», и план, построенный по таким оценкам, гоняет его по $5$ раз. EM и явная формула совпадают на всех тестах с $n_+\ge1$ (расхождение не больше $10^{-9}$).

\begin{center}\small
\begin{tabular}{@{}cccccccccc@{}}
\toprule
тест & $n_+$ & $S$ & $\hat p$ & $\hat s$ & $\hat p_{\rm EM}$ & $\hat s_{\rm EM}$ & итераций & $p$ & $s$ \\
\midrule
31 & 4 & 9 & 0.061 & 0.421 & 0.061 & 0.421 & 17 & 0.030 & 0.399 \\
36 & 1 & 5 & 0.014 & 1.000 & 0.014 & 1.000 & 2 & 0.001 & 0.890 \\
0 & 0 & 0 & 0.000 & --- & 0.000 & --- & 0 & 0.043 & 0.479 \\
8 & 8 & 21 & 0.118 & 0.510 & 0.118 & 0.510 & 12 & 0.131 & 0.516 \\
\bottomrule
\end{tabular}
\end{center}

Последние два столбца --- истинные параметры генератора (команды их не видят). Видно главное: при двух-трёх событиях на тест разделить $p$ и $s$ по отдельному тесту нельзя --- оценки уходят на границу или промахиваются в разы, хотя произведение $\hat p\hat s=S/(rN)$ оценивается честно. Лечится это только общей информацией по тестам: общая $s$ в группе (EM из замечания выше), априорное распределение на $s$ по всем сорока тестам (эмпирический Байес) или подстановка общей $\hat s$ тестам без повторных красных, как сделано в эталонном решении методом моментов.

\end{document}
